\section{向量场、曲线积分与势函数}
\label{sec:02-Calculus/06-Multivariable-Calculus/08}

\subsection{力沿弯路的功——一个自然的求和问题}

力乘以位移，中学就学过。但那是直线、恒力。现在想象一个小船在河流中航行：水流的速度在每一点都不同，方向各异，船走的又是一条弯曲的航线。发动机克服水流做了多少功？

你需要把整条航线切成极短的小段。每一段近似为直线，局部流速近似恒定，小段的功就是 \(\mathbf F\cdot\Delta\mathbf r\)——力在位移方向上的投影乘以位移长度。把所有小段累加起来，取极限，就是这个量的精确定义。

这就引出了向量场的概念。向量场

\[
\mathbf F: D\subseteq\R^n\to\R^n
\]

给区域中每一点赋一个向量——可以是流速、电场力、重力，或者优化中每一点的负梯度方向。与之对照，标量场 \(f:D\to\R\) 只给每一点一个数（温度、高度、概率密度），你可以问``这里有多高''；向量场让你问``这里的箭头指向哪、有多大''。

\subsection{两类曲线积分，方向性质截然不同}

给定一条参数曲线 \(\mathbf r(t)\)，\(a\le t\le b\)，向量场沿曲线的积分定义为

\[
\int_C \mathbf F\cdot d\mathbf r
= \int_a^b \mathbf F(\mathbf r(t))\cdot \mathbf r'(t)\,dt.
\]

点积 \(\mathbf F(\mathbf r(t))\cdot\mathbf r'(t)\) 只保留了力在运动切向上的分量。垂直分量对做功没有贡献——你推箱子时，如果力完全垂直于运动方向，箱子不会沿那个方向移动，物理上不做功。

一个重要性质：如果反向走同一条路径，\(d\mathbf r\) 反号，积分反号。功依赖于方向。

还有另一类以``曲线积分''为名的运算：

\[
\int_C f\,ds
= \int_a^b f(\mathbf r(t))\,\|\mathbf r'(t)\|\,dt.
\]

这里被积的是标量场的值乘以弧长微元 \(\|\mathbf r'(t)\|dt\)。它累积线密度或沿线的总质量。方向反转时，\(ds\) 不反号（弧长永远是正的），所以这类积分不关心方向。

两类积分都被称为``曲线积分''，却因被积对象的不同而拥有完全相反的方向性质。使用时必须先判断：累的是向量投影（点积），还是标量乘以弧长。

\subsection{参数跑快跑慢不影响几何结果}

同一条有向曲线可以用不同的参数化来描述。从 \((0,0)\) 到 \((1,2)\) 的直线，写成 \(\mathbf r(t)=(t,2t),\;0\le t\le 1\) 可以，写成 \(\mathbf r(u)=(u^2,2u^2),\;0\le u\le 1\) 也可以——后者的参数点移动速度不同，但走过的几何轨迹完全一样。

两种写法算出的积分是否相同？\(\mathbf r'(t)dt\) 或 \(\|\mathbf r'(t)\|dt\) 中的导数恰好补偿了参数变化速率的差异：参数跑得快的地方，切向量大，但对应的时间区间短；参数跑得慢的地方则相反。只要参数变换是光滑单调的，积分值不变。

这意味着真正决定曲线积分的是几何路径本身和方向，不是参数的``钟表''走多快。

\subsection{一种特殊的场：功只取决于起点和终点}

有些力场具有一个令人意外的性质：物体从 A 走到 B，无论走直线还是绕远路，场做的功都完全一样。重力场就是如此——把物体从地面搬到十楼，走楼梯和坐电梯，克服重力做的功相同（忽略路径差异导致的微小高度变化）。静电场的库仑力也有同样性质。

这个性质的数学表述是：存在一个标量函数 \(\phi\)，使得

\[
\mathbf F = \nabla\phi.
\]

一旦存在这样的 \(\phi\)（称为势函数），链式法则立即给出

\[
\frac{d}{dt}\phi(\mathbf r(t))
= \nabla\phi(\mathbf r(t))\cdot \mathbf r'(t)
= \mathbf F(\mathbf r(t))\cdot \mathbf r'(t).
\]

也就是说，沿曲线的被积函数恰好是 \(\phi\) 沿路径的全导数。于是积分退化为端点值之差：

\[
\int_C \nabla\phi\cdot d\mathbf r
= \phi(\mathbf r(b)) - \phi(\mathbf r(a)).
\]

这就是线积分的 Newton--Leibniz 定理——单变量微积分基本定理在曲线上的自然推广。闭合曲线的积分自动为零，因为起点和终点重合。具有势函数的向量场称为保守场。

举个例子：\(\mathbf F=(2x,2y)=\nabla(x^2+y^2)\)。从 \((0,0)\) 到 \((1,2)\)，沿任何路径的积分都是

\[
\phi(1,2)-\phi(0,0) = 1^2+2^2=5.
\]

这就是势函数的威力：把整条弯路上的累积计算压缩成两个点的函数值之差。

\subsection{怎么判断一个场能不能写成梯度？}

如果 \(\mathbf F=(P,Q)\) 是保守场，即 \(\mathbf F=\nabla\phi\)，那么 \(P=\phi_x\)，\(Q=\phi_y\)。在偏导数连续的条件下，混合偏导相等（Clairaut 定理）给出

\[
P_y = \phi_{xy} = \phi_{yx} = Q_x.
\]

所以 \(P_y=Q_x\) 是保守场的必要条件。在三维中，这个条件推广为旋度为零：

\[
\nabla\times\mathbf F = \mathbf 0.
\]

旋度的严格定义在下一节展开，此处先把它理解为一个衡量场``局部旋转程度''的量——保守场没有``漩涡''。

\subsection{导数条件够了？还差一个拓扑门槛}

\(P_y=Q_x\) 或旋度为零是否能保证全局存在势函数？不能。经典的翻身案例是平面上挖去原点的场：

\[
\mathbf F(x,y) = \left(-\frac{y}{x^2+y^2},\; \frac{x}{x^2+y^2}\right).
\]

它在定义域（去掉原点的全平面）内旋度处处为零，却沿绕原点逆时针的单位圆积分得到 \(2\pi\)，不为零，因而不保守。

出了问题的地方在哪？原点处的洞。你可以在定义域内画出围绕原点的闭合曲线，却无法在定义域内把它连续收缩成一个点——原点不让你穿过去。这个洞``锁住''了全局环流，局部导数条件检测不到它。

如果定义域满足``任何闭合曲线都能在域内连续收缩成一点''（这样的区域称为单连通），那么平面中的 \(P_y=Q_x\) 或三维中的 \(\nabla\times\mathbf F=\mathbf0\) 加上光滑性条件，就能保证势函数存在。简单连通域（例如整个平面、球体内部、凸区域）上，旋度为零等价于保守性。

\subsection{找势函数的手工活：积分``常数''其实是另一个变量的函数}

知道一个场是保守场后，怎么把势函数找出来？设 \(\phi_x=P(x,y)\)。对 \(x\) 积分：

\[
\phi(x,y) = \int P(x,y)\,dx + C(y).
\]

关键在这里：对 \(x\) 不定积分时，\(y\) 是被冻结的，所以积分``常数''可以是 \(y\) 的任意函数。很多人在这一步只写一个 \(+C\) 常数就停了，这会把势函数丢掉一部分。

有了带 \(C(y)\) 的表达式后，再对 \(y\) 求偏导，令其等于 \(Q(x,y)\)：

\[
\phi_y = \frac{\partial}{\partial y}\int P(x,y)\,dx + C'(y) = Q(x,y).
\]

从中解出 \(C'(y)\)，再积分得到 \(C(y)\)。三维同理：先对一个变量积分，积分``常数''依赖其余变量，再依次用其他偏导条件定出。

整个过程是反向的梯度运算：梯度把势函数拆成各分量，找势函数就是把分量重新拼回去，顺带恢复被偏导数抹掉的``跨变量''信息。
